\chapter{System Architecture and Software Design}
\label{ch:architecture}

\section{Overall Architecture}
The application follows a three-tier architecture, but unlike a classic
server-rendered PHP application, the frontend is a React single-page
application that communicates with the backend purely through a JSON API,
rather than through full-page navigations. Figure~\ref{fig:architecture}
shows the overall structure.

\begin{figure}[h]
\centering
\begin{tikzpicture}[node distance=1.4cm, every node/.style={font=\small}]
    \tikzstyle{tier} = [rectangle, draw=black, thick, rounded corners,
        fill=orange!10, text width=8cm, align=center, minimum height=1.1cm]
    \node (client) [tier] {React SPA (Vite) \\ Context, Hooks, Components, Services};
    \node (api) [tier, below=of client] {Fetch API calls (JSON, \texttt{credentials: include})};
    \node (backend) [tier, below=of api] {PHP Endpoints (Procedural, PDO)};
    \node (db) [tier, below=of backend] {MySQL Database (\texttt{users}, \texttt{game\_results})};

    \draw [-{Latex[length=2mm]}, thick] (client) -- (api);
    \draw [-{Latex[length=2mm]}, thick] (api) -- (backend);
    \draw [-{Latex[length=2mm]}, thick] (backend) -- (db);
    \draw [-{Latex[length=2mm]}, thick] (db) -- (backend);
    \draw [-{Latex[length=2mm]}, thick] (backend) -- (api);
    \draw [-{Latex[length=2mm]}, thick] (api) -- (client);
\end{tikzpicture}
\caption{Overall three-tier architecture of the application.}
\label{fig:architecture}
\end{figure}

\section{Frontend Architecture}
The frontend is organised into four layers:
\begin{itemize}[nosep]
    \item \textbf{Context providers} -- \texttt{AuthContext} holds the
    current user and authentication status; \texttt{ThemeContext} holds the
    light/dark theme preference; a \texttt{ToastContext} manages transient
    notification messages.
    \item \textbf{Custom hooks} -- \texttt{useMinesweeper} is the central
    state machine for a single game session, wrapping the pure functions
    from Chapter~\ref{ch:game-design} with React state and effects (the
    countdown timer, hint highlight timeout, and automatic score
    submission).
    \item \textbf{Components / Pages} -- page-level components
    (Login, Register, Dashboard, Game, Leaderboard, History) composed from
    smaller presentational components (the board grid, a cell, the
    stats cards).
    \item \textbf{Services} -- \texttt{api.js} wraps the browser
    \texttt{fetch} API with a consistent JSON request/response contract;
    \texttt{authService.js} and \texttt{gameService.js} expose
    typed functions (e.g.\ \texttt{login}, \texttt{register},
    \texttt{saveScore}, \texttt{getLeaderboard}) built on top of it.
\end{itemize}

\begin{figure}[h]
\centering
\begin{tikzpicture}[node distance=1.2cm, every node/.style={font=\footnotesize}]
    \tikzstyle{box} = [rectangle, draw=black, rounded corners, fill=blue!8,
        text width=3.6cm, align=center, minimum height=0.9cm]
    \node (ctx) [box] {Context Providers \\ (Auth, Theme, Toast)};
    \node (hooks) [box, below=of ctx] {Custom Hooks \\ (useMinesweeper)};
    \node (comp) [box, below=of hooks] {Pages / Components};
    \node (svc) [box, below=of comp] {Services (api.js, authService, gameService)};
    \node (backend) [box, below=of svc, fill=orange!10] {PHP Backend};

    \draw [-{Latex[length=2mm]}, thick] (ctx) -- (hooks);
    \draw [-{Latex[length=2mm]}, thick] (hooks) -- (comp);
    \draw [-{Latex[length=2mm]}, thick] (comp) -- (svc);
    \draw [-{Latex[length=2mm]}, thick] (svc) -- (backend);
\end{tikzpicture}
\caption{Frontend layered data flow.}
\label{fig:frontend-layers}
\end{figure}

\section{Backend Architecture}
The backend is written in procedural PHP, organised as one script per
endpoint (e.g.\ \texttt{register.php}, \texttt{login.php},
\texttt{save-score.php}, \texttt{leaderboard.php}, \texttt{history.php},
\texttt{stats.php}), plus shared includes for the database connection and
session/authentication helpers. Every protected endpoint begins by calling a
shared \texttt{require\_login()} helper, which checks the PHP session for a
logged-in user ID and terminates the request with an error response if it
is absent.

\section{Client-Side Route Guarding vs.\ Server-Side Authorization}
A React \texttt{ProtectedRoute} component wraps pages such as the dashboard
and the game page, redirecting to the login page if \texttt{AuthContext}
does not currently hold an authenticated user. It is important to note that
this client-side check is a \emph{user-experience} convenience only: it
prevents an unauthenticated user from momentarily seeing a protected page's
markup, but it grants no actual access to data. The real security boundary
is enforced entirely server-side, by \texttt{require\_login()} checking the
PHP session on every protected endpoint (Section~\ref{sec:session-security}
in Chapter~\ref{ch:security}). A user could bypass \texttt{ProtectedRoute}
entirely (e.g.\ by disabling JavaScript route checks) and would still be
unable to retrieve another user's data, because the backend independently
authenticates every request.

\section{API Communication}
All requests from the frontend to the backend are issued with
\texttt{credentials: "include"}, so that the PHP session cookie is sent with
cross-origin requests. Responses are JSON objects with a consistent shape
(a success flag and either a data payload or an error message), which the
\texttt{api.js} wrapper parses uniformly for every service call.
